\section{Molekulyar kinitik nazaryasi(MKN) asosiy tushunchalari}


\begin{enumerate}

    \item Molekular-kinetik nazariyaning asosiy qoidalari qaysi javobda nisbatan to‘liq bayon qilingan? 
    \begin{itemize}
        \item[A)] har qanday jism molekulalardan tuzilgan, ular betartib harakatda bo‘ladi, ular orasida o‘zaro ta’sir kuchlari mavjud.
        \item[B)] modda mayda zarralardan tashkil topgan va ular orasida o‘zaro ta’sir kuchlari mavjud.
        \item[C)] modda mayda bo‘laklardan tuziladi va ular fazoni uzluksiz to‘ldiradi.
        \item[D)] modda elementar zarralardan tuziladi va ular bir-biriga tinimsiz aylanib turadi.
     \end{itemize}

    \item 1 mol BXS da qaysi kattalikning birligidir?
    \begin{itemize}
        \item[A)] massa.
        \item[B)] modda miqdori.
        \item[C)] molyar massa.
        \item[D)] hajm.
       
    \end{itemize}

    \item Modda miqdorining birligi nima?
    \begin{itemize}
        \item[A)] mol.
        \item[B)] kg.
        \item[C)] kg/mol.
        \item[D)] m.a.b.
        
    \end{itemize}

    \item Avagadro soni deb qanday fizik kattalikka aytiladi? 
    \begin{itemize}
        \item[A)] 12 g ugleroddagi atomlar soniga.
        \item[B)] bitta uglerod atomi massasining 1/12 qismiga.
        \item[C)] 1 mol modda massasi.
        \item[D)] bitta atom massasining uglerod atomi massasining 1/12 qismiga nisbatiga.
        
    \end{itemize}

    \item Misning nisbiy atom massasi 64 m.a.b. Bitta mis atomining massasini toping(kg). $1 \text{ m.a.b} = 1,67 \cdot 10^{-27} \text{ kg}$. [2, 8]
    \begin{itemize}
        \item[A)] $1,2 \cdot 10^{-26}$.
        \item[B)] $1,5 \cdot 10^{-26}$.
        \item[C)] $2,03 \cdot 10^{-25}$.
        \item[D)] $1,07 \cdot 10^{-25}$.
        
    \end{itemize}

    \item Bitta oltin atomining massasi $3,29 \cdot 10^{-25} \text{ kg}$. Uning nisbiy atom massasini toping (m.a.b).
    \begin{itemize}
        \item[A)] 32.
        \item[B)] 108.
        \item[C)] 197.
        \item[D)] 29.
        
    \end{itemize}

    \item Molekulalar soni $1,806 \cdot 10^{23}$ ta bo‘lgan moddaning miqdorini toping (mol). Avogadro soni $N_A = 6,02 \cdot 10^{23} \text{ mol}^{-1}$.
    \begin{itemize}
        \item[A)] 0,8.
        \item[B)] 0,3.
        \item[C)] 0,2.
        \item[D)] 0,5.
        
    \end{itemize}

    \item Molekulalar soni $15,05 \cdot 10^{24}$ ta bo‘lgan moddaning miqdorini toping (mol). $N_A = 6,02 \cdot 10^{23} \text{ mol}^{-1}$.
    \begin{itemize}
        \item[A)] 15.
        \item[B)] 18.
        \item[C)] 25.
        \item[D)] 20.
        
    \end{itemize}

    \item 39,4 kg oltindagi modda miqdorini aniqlang(mol). $\mu = 197 \text{ g/mol}$.
    \begin{itemize}
        \item[A)] 400.
        \item[B)] 200.
        \item[C)] 100.
        \item[D)] 150.
        
    \end{itemize}

    \item 9 g suvda necha mol modda bor? $\mu = 18 \text{ g/mol}$. 
    \begin{itemize}
        \item[A)] 0,45.
        \item[B)] 0,5.
        \item[C)] 0,9.
        \item[D)] 4,5.
        
    \end{itemize}

    \item 28,8 g kislorod moddasining miqdorini (mol) aniqlang. $\mu = 32 \cdot 10^{-3} \text{ kg/mol}$.
    \begin{itemize}
        \item[A)] 0,5.
        \item[B)] 1.
        \item[C)] 2.
        \item[E)] 0,9.
    \end{itemize}

    \item 5 mol geliyning massasi qanday? $\mu = 4 \text{ g/mol}$.
    \begin{itemize}
        \item[A)] 0,1 kg.
        \item[B)] 0,2 kg.
        \item[C)] 10 g.
        \item[D)] 20 g.
        
    \end{itemize}

    \item 450 mol kislorodning massasi qanday (kg)? $\mu = 32 \text{ g/mol}$.
    \begin{itemize}
        \item[A)] 14,4.
        \item[B)] 4,5.
        \item[C)] 7,2.
        \item[D)] 9.
        
    \end{itemize}

    \item 5 mol kislorodning massasini toping. Kislorodning molyar massasi $32 \cdot 10^{-3} \text{ kg/mol}$.
    \begin{itemize}
        \item[A)] 320 g.
        \item[B)] 32 g.
        \item[C)] 1,6 kg.
        \item[D)] 160 g.
        \
    \end{itemize}

    \item Modda miqdorlari teng bo‘lgan qalay va nikel quymalarning massalarini taqqoslang. Qalayning molyar massasi $\mu_q = 118,7 \text{ g/mol}$, nikelniki esa $\mu_n = 58,7 \text{ g/mol}$.
    \begin{itemize}
        \item[A)] $m_q \approx 4m_n$.
        \item[B)] $m_n \approx 2m_q$.
        \item[C)] $m_q \approx 2m_n$.
        \item[D)] $m_n \approx 4m_q$.
        
    \end{itemize}

    \item Butan ($C_4H_{10}$) ning molyar massasini toping (kg/mol). $Ar(H)=1 \text{ u}, Ar(C)=12 \text{ u}, 1 \text{ u} = 1,66 \cdot 10^{-27} \text{ kg}$.
    \begin{itemize}
        \item[A)] $17 \cdot 10^{-3}$.
        \item[B)] $41 \cdot 10^{-3}$.
        \item[C)] $80 \cdot 10^{-3}$.
        \item[D)] $58 \cdot 10^{-3}$.
        
    \end{itemize}

    \item Azot kislota ($HNO_3$) ning molyar massasini toping (kg/mol). $Ar(H)=1 \text{ u}, Ar(N)=14 \text{ u}, Ar(O)=16 \text{ u}, 1 \text{ u} = 1,66 \cdot 10^{-27} \text{ kg}$.
    \begin{itemize}
        \item[A)] $56 \cdot 10^{-3}$.
        \item[B)] $63 \cdot 10^{-3}$.
        \item[C)] $31 \cdot 10^{-3}$.
        \item[D)] $15 \cdot 10^{-3}$.
        
    \end{itemize}

    \item Mis kuporosi ($CuSO_4$) ning uch moli massasini aniqlang (g). $Ar(O)=16 \text{ u}, Ar(Cu)=64 \text{ u}, Ar(S)=32 \text{ u}, 1 \text{ m.a.b.} = 1,66 \cdot 10^{-27} \text{ kg}, N_A = 6 \cdot 10^{23} \text{ mol}^{-1}$.
    \begin{itemize}
        \item[A)] 160.
        \item[B)] 540.
        \item[C)] 480.
        \item[D)] 240.
       
    \end{itemize}

    \item Gazning bitta molekulasining massasi $4,8 \cdot 10^{-26} \text{ kg}$ ga teng. Shu gazning molyar massasi qanday (g/mol). $N_A = 6,02 \cdot 10^{23} \text{ mol}^{-1}$.
    \begin{itemize}
        \item[A)] 8.
        \item[B)] 32.
        \item[C)] 12.
        \item[D)] 2.
       
    \end{itemize}

    \item Moddaning molyar massasi 36 g/mol bo‘lsa, bitta molekulasining massasi necha kilogramm bo‘ladi? $N_A = 6 \cdot 10^{23} \text{ mol}^{-1}$.
    \begin{itemize}
        \item[A)] $6 \cdot 10^{-23}$.
        \item[B)] $6 \cdot 10^{-26}$.
        \item[C)] $216 \cdot 10^{-20}$.
        \item[D)] $36 \cdot 10^{-23}$.
        
    \end{itemize}

    \item Azotning molyar massasi 0,028 kg/mol. Azot molekulasining massasi qanday (kg)? $N_A = 6 \cdot 10^{23} \text{ mol}^{-1}$.
    \begin{itemize}
        \item[A)] $5,3 \cdot 10^{-23}$.
        \item[B)] $4,7 \cdot 10^{-26}$.
        \item[C)] $5,3 \cdot 10^{-27}$.
        \item[D)] $8,6 \cdot 10^{-22}$.
       
    \end{itemize}

    \item Piyolaga quyilgan 180 g suv bir haftada to‘la bug‘lanib ketadi. Piyoladan bir sekundda o‘rta hisobda nechta molekula bug‘lanib turganligini baholang. $N_A = 6 \cdot 10^{23} \text{ mol}^{-1}, \mu = 18 \text{ g/mol}$.
    \begin{itemize}
        \item[A)] $10^{22}$.
        \item[B)] $10^{19}$.
        \item[C)] $10^{15}$.
        \item[D)] $10^9$.
        
    \end{itemize}

    \item Idishdagi gaz massasining 84 foizini azot, 16 foizini kislorod tashkil qilsa, bitta kislorod molekulasiga nechta azot molekulasi to‘g‘ri keladi? $M_A = 28 \text{ g/mol}, M_K = 32 \text{ g/mol}$.
    \begin{itemize}
        \item[A)] 4.
        \item[B)] 3.
        \item[C)] 2.
        \item[D)] 6.
       
    \end{itemize}

    \item Sirt yuzasi $20 cm^2$ bo‘lgan qurilmaning sirtiga $1m$ qalinlikda oltin qatlami qoplandi. Qoplamda nechta oltin atomi bor? $N_A = 6 \cdot 10^{23} \text{ mol}^{-1}$, oltinning atom massasi 197 m.a.b., zichligi esa $19,3 \text{ g/cm}^3$.
    \begin{itemize}
        \item[A)] $1,18 \cdot 10^{20}$.
        \item[B)] $1,5 \cdot 10^{20}$.
        \item[C)] $2 \cdot 10^{20}$.
        \item[D)] $1,13 \cdot 10^{20}$.
        
    \end{itemize}

    \item Bitta molekula massasi $6 \cdot 10^{-26} \text{ kg}$ va molekulalar konsentratsiyasi $2,5 \cdot 10^{29} \text{ m}^{-3}$ bo‘lgan metallning zichligini toping (kg/m$^3$).
    \begin{itemize}
        \item[A)] 1200.
        \item[B)] 6000.
        \item[C)] 15000.
        \item[D)] 2000.
       
    \end{itemize}

    \item Diffuziya qanday moddalarda kuzatiladi?
    \begin{itemize}
        \item[A)] Faqat gazlarda.
        \item[B)] Gazlar, suyuq va qattiq jismlarda.
        \item[C)] Faqat suyuqlik va qattiq jismlarda.
        \item[D)] Faqat gaz va suyuqliklarda.
       
    \end{itemize}

    \item Broun harakati deyilganda, nima tushuniladi? 
    \begin{itemize}
        \item[A)] suyuqlik molekulalarining issiqlik harakati.
        \item[B)] suyuqlikka solingan qattiq jism zarralarining suyuqlik molekulalari ta’siridagi harakati.
        \item[C)] suyuqlik molekulalarining unga solingan mayda zarralar ta’siridagi harakati.
        \item[D)] suyuqlikka solingan qattiq jism zarralarining o‘zaro ta’sirlashuvlari natijasida yuzaga keladigan harakati.
    \end{itemize}

    \item Idishdagi kislorod molekulalarining konsentratsiyasi $n = 6 \cdot 10^{25} \text{ m}^{-3}$ bo‘lsa, gaz zichligi qanday (kg/m$^3$)? $N_A = 6 \cdot 10^{23} \text{ mol}^{-1}, \mu = 32 \cdot 10^{-3} \text{ kg/mol}$.
    \begin{itemize}
        \item[A)] 0,6.
        \item[B)] 1,2.
        \item[C)] 1,6.
        \item[D)] 3,2.
        
    \end{itemize}

    \item Massasi 320 g bo‘lgan kislorod $0,02 \text{ m}^3$ hajmni egallab turibdi. Molekulalar konsentratsiyasini toping (m$^{-3}$). $\mu = 32 \text{ g/mol}, N_A = 6 \cdot 10^{23} \text{ mol}^{-1}$.
    \begin{itemize}
        \item[A)] $3,01 \cdot 10^{27}$.
        \item[B)] $3,01 \cdot 10^{26}$.
        \item[C)] $6,02 \cdot 10^{21}$.
        \item[D)] $6,02 \cdot 10^{27}$.
       
    \end{itemize}

\end{enumerate}











\begin{table}[!ht]
\centering
\begin{tabular}{|c|c|c|c|c|c|c|c|c|c|c|}
\hline
1a & 2b & 3a & 4a & 5d & 6c & 7b & 8c & 9b & 10b & 11d \\ \hline
12d & 13a & 14d & 15c & 16d & 17b & 18c & 19a & 20b & 21b & 22b \\ \hline
23d & 24a & 25d & 26b & 27b & 28d & 29b &  &  &  &  \\ \hline
\end{tabular}
\end{table}